% =============================================================================
% Project : THE_PHILOSOPHY_OF_VIABILITY
% Book    : Book I — Ontology of Viability
% Book ID : VS.B01
%
% File    : 4_ASSETS/tables/table_01_forms_of_invariance.tex
%
% Section / Chapter / Asset:
%   Table TL01 — Analytical Forms of Invariance
%
% Asset ID:
%   VS.B01.P01.C02.S01.SS05.P001.TL01.forms_of_invariance
%
% Purpose :
%   To distinguish analytically different respects in which
%   determinations, relations, structures, compositions, and
%   configurations may be preserved under transformation.
%
% Conceptual Policy:
%   1. The categories identify analytical respects, not necessarily
%      mutually exclusive ontological kinds.
%   2. Qualitative and quantitative invariance are treated as
%      subtypes of determinational invariance.
%   3. Structural, compositional, and configurational invariance
%      are distinguished without presupposing absolute independence.
%   4. No form of invariance is presumed to apply universally.
%   5. The classification remains provisional and open to refinement.
%
% =============================================================================

\begin{center}
\small
\renewcommand{\arraystretch}{1.28}
\setlength{\tabcolsep}{4pt}

\begin{longtable}{
    >{\raggedright\arraybackslash}p{2.6cm}
    >{\raggedright\arraybackslash}p{3.0cm}
    >{\raggedright\arraybackslash}p{3.2cm}
    >{\raggedright\arraybackslash}p{4.4cm}
}

\caption{Analytical Forms and Respects of Invariance}
\label{tab:forms-of-invariance} \\

\toprule
\textbf{Analytical Respect}
&
\textbf{What Is Claimed to Be Preserved}
&
\textbf{Relevant Transformation}
&
\textbf{Ontological Qualification}
\\
\midrule
\endfirsthead

\toprule
\textbf{Analytical Respect}
&
\textbf{What Is Claimed to Be Preserved}
&
\textbf{Relevant Transformation}
&
\textbf{Ontological Qualification}
\\
\midrule
\endhead

\midrule
\multicolumn{4}{r}{\footnotesize\textit{Continued on next page}}\\
\endfoot

\bottomrule
\endlastfoot

Determinational Invariance
&
A specified determination or feature of an existent.
&
Transformations that may alter other determinations
while leaving the specified feature preserved.
&
The preserved determination does not thereby
exhaust the constitution of the existent or
establish its numerical identity.
\\

\addlinespace

\quad Qualitative Invariance
&
A specified qualitative determination.
&
Transformations that may alter other qualitative
or quantitative aspects.
&
The preservation of a quality must be established
according to an explicit criterion; resemblance
alone is insufficient.
\\

\addlinespace

\quad Quantitative Invariance
&
A specified magnitude, number, ratio, or other
quantitative determination.
&
Transformations that may alter the configuration
or other features of the existent.
&
Quantitative preservation does not establish
qualitative, structural, compositional, or
numerical identity.
\\

\midrule

Relational Invariance
&
A specified relation among relata.
&
Transformations affecting the relata, their
determinations, or their configurations.
&
A relation may remain preserved despite changes
in the relata. The criterion of relational
preservation must be specified.
\\

\addlinespace

Structural Invariance
&
A specified organizational pattern or structural
relation.
&
Transformations affecting components or particular
determinations while potentially preserving
an organizational arrangement.
&
The preservation of structure does not necessarily
entail the preservation of the same components,
composition, or numerical identity.
\\

\addlinespace

Compositional Invariance
&
A specified aspect of the composition of an existent,
including its constituents or compositional relations.
&
Transformations affecting components, their
inclusion, replacement, arrangement, or
compositional relations.
&
The relevant respect must clarify whether
preservation concerns constituents, compositional
relations, or another aspect of constitution.
\\

\addlinespace

Configurational Invariance
&
A specified configuration or arrangement of
determinations.
&
Transformations that may alter underlying elements
or other aspects while potentially preserving
the relevant arrangement.
&
The criterion of configuration must be explicit.
Similarity between configurations does not by
itself establish invariance.
\\

\end{longtable}

\vspace{0.5em}

\begin{minipage}{0.95\textwidth}
\footnotesize

\textit{Note:}
The categories distinguish different analytical
respects in which preservation may be considered.
They are not asserted to be mutually exclusive,
exhaustive, or ontologically independent.

Qualitative and quantitative invariance are presented
as subtypes of determinational invariance. Relational
and structural respects may jointly characterize
a case of preservation without thereby constituting
a separate fundamental kind. Compositional and
configurational invariance likewise require
specification of the particular aspect claimed
to remain preserved.

A single transformation may preserve several
determinations or relations while altering others.
The table does not establish that any particular
invariant exists in every existent or under
every transformation.

\end{minipage}
\end{center}

% =============================================================================
% END OF TABLE TL01
% =============================================================================